\documentclass[11pt]{article}
\usepackage[margin=1in]{geometry}
\usepackage{amsmath,amssymb,amsthm}
\usepackage{booktabs}
\usepackage{hyperref}

\newtheorem{proposition}{Proposition}
\newtheorem{definition}{Definition}

\title{Predictive validation of a transfer-function model:\\
       the out-of-sample forecast diagnosis of \texttt{drtran}}
\author{drtran / ATSW suite}
\date{}

\begin{document}
\maketitle

\begin{abstract}
A multivariate transfer-function model is worth building \emph{only if it
forecasts better than the univariate model it extends}. This note fixes the
mathematics of the out-of-sample comparison that \texttt{drtran} uses to decide
this: a fixed-parameter recursive forecast over balanced origins, the
root-mean-square error by horizon, the ratio between two models, and the
Diebold--Mariano test with its heteroskedasticity-and-autocorrelation-consistent
(HAC) variance and Harvey--Leybourne--Newbold (HLN) small-sample correction. It
also makes explicit the overlap structure of multi-step forecasts and the
resulting effective sample size, which governs how much the long-horizon test
can be trusted.
\end{abstract}

\section{The instrumentalist criterion}

Let $A$ be a univariate model for an output series $y_t$ and $B$ a multivariate
model that adds one or more inputs (a transfer function, $y_t=\nu(B)x_t+N_t$).
$B$ nests $A$: with $\nu\equiv 0$, $B$ collapses to $A$. In sample, $B$ can never
fit worse than $A$ (the likelihood is monotone in the free parameters); a
significant $\hat\omega$ or a higher log-likelihood is therefore \emph{necessary
but not sufficient} to justify $B$.

Following Friedman's instrumentalism, the model is judged by what it predicts.
The operational criterion is:
\[
\boxed{\;B \text{ is justified over } A \iff B \text{ forecasts } y \text{ better
out of sample than } A.\;}
\]
Everything below is the machinery to make ``better out of sample'' precise and
testable.

\section{The fixed-parameter recursive forecast}

Let $y_1,\dots,y_n$ be the observed output (and $x_1,\dots,x_n$ the inputs).
Split the sample at $E$ (the \emph{estimation window}, \texttt{-estwin E}): the
training span is $1,\dots,E$ and the evaluation span is $E+1,\dots,n$.

\begin{definition}[Fixed-parameter forecast from an origin]
Estimate the model once on $1,\dots,E$, obtaining $\hat\theta_E$. For a
\emph{forecast origin} $e\ge E$, condition on the data through $e$ and produce the
$h$-step-ahead point forecast of the output,
\[
\hat y_{e+h\mid e}(\hat\theta_E), \qquad h=1,\dots,H,
\]
holding $\hat\theta_E$ \emph{fixed} (no re-estimation as $e$ advances). For a
transfer model, $\hat y_{e+h\mid e}$ requires the input forecast
$\hat x_{e+j\mid e}$, so the input's own model is projected forward too.
\end{definition}

The parameters come from the training window $E$; the origin $e$ is independent of
$E$ (\texttt{drtran} exposes them as \texttt{-estwin E} and \texttt{-O e}). The SPS
real-time use sets $e=n$ (the current data end); the evaluation use rolls $e$ over
the hold-out (below).

\section{Balanced origins and the error panel}

\begin{definition}[Balanced origins]
For a fixed horizon $H$, roll the origin over
\[
e = E,\,E+1,\,\dots,\,n-H,
\]
so that \emph{every} origin has a full set of $H$ realised values
$y_{e+1},\dots,y_{e+H}$ to be scored against. There are $N=n-H-E+1$ origins, and
each contributes exactly one error at each horizon $h\le H$. Hence every horizon is
scored on the same number of errors $N$ (no horizon is scored on fewer, noisier
points).
\end{definition}

Define the $h$-step forecast error of model $M\in\{A,B\}$ from origin $e$:
\[
\varepsilon^{M}_{e,h} \;=\; \hat y^{M}_{e+h\mid e}(\hat\theta_E) \;-\; y_{e+h}.
\]
This is the \emph{error panel} $\{\varepsilon^M_{e,h}\}$, indexed by origin $e$ and
horizon $h$ --- exactly the CSV that \texttt{drtran -C} writes
(\texttt{origin,horizon,actual,forecast,error}).

\section{Accuracy by horizon and the ratio}

For each horizon $h$, aggregate over the $N$ origins:
\[
\mathrm{RMSE}^{M}(h)=\sqrt{\frac1N\sum_{e=E}^{\,n-H}\!\big(\varepsilon^{M}_{e,h}\big)^2},
\qquad
\mathrm{MAE}^{M}(h)=\frac1N\sum_{e}\big|\varepsilon^{M}_{e,h}\big|.
\]
The head-to-head summary is the ratio
\[
r(h)=\frac{\mathrm{RMSE}^{B}(h)}{\mathrm{RMSE}^{A}(h)},
\qquad r(h)<1 \iff B \text{ more accurate at horizon } h.
\]

\section{The Diebold--Mariano test}

The ratio is a point estimate; whether $B$ beats $A$ \emph{significantly} at
horizon $h$ is a hypothesis test on the loss differential. With squared-error loss,
define, for each origin $e$,
\[
d_{e,h} \;=\; \big(\varepsilon^{A}_{e,h}\big)^2 - \big(\varepsilon^{B}_{e,h}\big)^2,
\qquad
\bar d_h=\frac1N\sum_e d_{e,h}.
\]
$\bar d_h>0$ means $B$ has the smaller mean squared error. Diebold and Mariano
(1995) test $H_0:\ \mathbb E[d_{e,h}]=0$ (equal predictive accuracy) against
$H_1:\ \mathbb E[d_{e,h}]\neq 0$.

\paragraph{Why a HAC variance is unavoidable.}
The $h$-step forecasts from consecutive origins overlap.

\begin{proposition}[Overlap induces $\mathrm{MA}(h-1)$ dependence]
Two $h$-step-ahead forecasts made from origins $e$ and $e+k$ target the windows
$[e+1,e+h]$ and $[e+k+1,e+k+h]$, which \emph{overlap} whenever $k<h$. Under
correct specification the $h$-step error is a moving average of the $h$ innovations
entering the forecast window, so $\{d_{e,h}\}_e$ is serially correlated up to lag
$h-1$ and uncorrelated beyond. Hence a consistent estimate of
$\mathrm{Var}(\bar d_h)$ must sum autocovariances to lag $h-1$.
\end{proposition}

Accordingly, use the rectangular (truncated) HAC long-run variance to lag $h-1$,
\[
\widehat{\mathrm{LRV}}_h
=\hat\gamma_0+2\sum_{k=1}^{h-1}\hat\gamma_k,
\qquad
\hat\gamma_k=\frac1N\sum_{e}\big(d_{e,h}-\bar d_h\big)\big(d_{e+k,h}-\bar d_h\big).
\]
The raw statistic is $\mathrm{DM}_h=\bar d_h\big/\sqrt{\widehat{\mathrm{LRV}}_h/N}$.

\paragraph{Small-sample (HLN) correction.}
Harvey, Leybourne and Newbold (1997) show $\mathrm{DM}_h$ is oversized in small
samples at long horizons and propose the correction factor
\[
c_{N,h}=\sqrt{\frac{N+1-2h+h(h-1)/N}{N}},
\qquad
\mathrm{DM}^{\*}_h=c_{N,h}\,\mathrm{DM}_h,
\]
compared against the Student-$t$ distribution with $N-1$ degrees of freedom:
\[
p_h = 2\,\Pr\!\big(t_{N-1} > |\mathrm{DM}^{\*}_h|\big).
\]
(If $\widehat{\mathrm{LRV}}_h\le 0$ the estimator falls back to the i.i.d.
variance $\hat\gamma_0$.) This is the exact statistic implemented in
\texttt{forecast\_compare.py::dm}.

\section{Overlap and the effective sample size}

The RMSE at horizon $h$ is a clean point estimate: it averages $N$ errors, the
same $N$ for every $h$ (balanced origins). The \emph{test}, however, cannot treat
those $N$ errors as independent, because they overlap. A set of \emph{non-overlapping}
$h$-step forecasts is obtained by spacing origins $h$ apart, giving
\[
N_{\mathrm{eff}}(h) \;\approx\; \frac{N}{h}
\]
effectively independent blocks. Thus:
\begin{itemize}
\item at $h=1$, $N_{\mathrm{eff}}=N$ --- the test is at full strength;
\item at long horizons, $N_{\mathrm{eff}}\to N/h$ becomes tiny (e.g.\ $N=37,\ h=24
      \Rightarrow N_{\mathrm{eff}}\approx1.5$).
\end{itemize}
The HAC-to-$(h-1)$ and HLN correction are precisely what keep the test honest under
this overlap, but no correction manufactures independent data. \textbf{Reading
rule:} the \emph{ratio} $r(h)$ and its \emph{sign} are trustworthy at all horizons;
the \emph{$p$-value} is strong at short horizons and fragile at long horizons, where
$N_{\mathrm{eff}}$ is small. \texttt{forecast\_compare.py} prints $N_{\mathrm{eff}}$
alongside $N$ for this reason.

\section{The verdict}

For horizons $h\in\{1,2,6,12,24\}$ (short to long), report
$\big(N,\ N_{\mathrm{eff}},\ \mathrm{RMSE}^A,\ \mathrm{RMSE}^B,\ r,\ \mathrm{DM}^{\*},\ p\big)$.
The instrumentalist verdict is:
\[
\text{keep } B \text{ over } A \iff r(h)<1 \text{ (materially, and with }
\mathrm{DM}^{\*}_h>0\text{) over the horizons of interest.}
\]
A model that fits better in sample but does not lower $r(h)$ below $1$ out of
sample fails the criterion. Because $r(h)$ can depend on the regime of the
evaluation span (the input may be predictable in calm periods and unpredictable in
crises), the diagnosis should be run over more than one hold-out period; a sign
change of $r-1$ across regimes is itself the finding.

\section*{References}
\begin{itemize}\itemsep0pt
\item Diebold, F.X.\ and Mariano, R.S.\ (1995). Comparing predictive accuracy.
      \emph{J.\ Business \& Economic Statistics} 13(3), 253--263.
\item Harvey, D., Leybourne, S.\ and Newbold, P.\ (1997). Testing the equality of
      prediction mean squared errors. \emph{Int.\ J.\ Forecasting} 13(2), 281--291.
\item Abraham, B.\ and Box, G.E.P.\ (1978). Deterministic and forecast-adaptive
      time-dependent models. \emph{Applied Statistics} 27(2), 120--130.
\item Friedman, M.\ (1953). The methodology of positive economics. In
      \emph{Essays in Positive Economics}, Univ.\ of Chicago Press.
\end{itemize}

\end{document}
